% Part: turing-machines
% Chapter: machines-computations
% Section: unary-numbers

\documentclass[../../../include/open-logic-section]{subfiles}

\begin{document}

\olfileid[pt]{tur}{mac}{una}
\olsection{Representação Unária de Números}

\begin{explain}
As máquinas de Turing trabalham sobre sequências de símbolos escritos em sua fita.
Dependendo do alfabeto que uma máquina de Turing usa, essas sequências de
símbolos podem representar diversas entradas e saídas.  De particular
interesse, é claro, estão as máquinas de Turing que computam
funções \emph{aritméticas}, isto é, funções de números naturais.  Uma
maneira simples de representar inteiros positivos é codificá-los como
sequências de um único símbolo~$\TMstroke$.  Se $n \in \Nat$, seja
$\TMstroke^n$ a sequência vazia se $n = 0$, e caso contrário a
sequência composta de exatamente $n$ $\TMstroke$.
\end{explain}

\begin{defn}[Computação]
Uma máquina de Turing~$M$ \emph{computa} a função $f\colon \Nat^k \to \Nat$ sse
$M$~para com a entrada
\[
\TMstroke^{n_1} \TMblank \TMstroke^{n_2} \TMblank \dots \TMblank \TMstroke^{n_k}
\]
produzindo a saída $\TMstroke^{f(n_1, \dots, n_k)}$.
\end{defn}

\begin{prob}
  Dê uma definição para quando uma máquina de Turing~$M$ computa a
  função $f\colon \Nat^k \to \Nat^m$.
\end{prob}

\begin{ex}\ollabel{ex:adder}
\emph{Adição:}
Vamos construir uma máquina que computa a função $f(n,m) = n + m$.
Isso requer uma máquina que comece com dois blocos de $\TMstroke$
de comprimento $n$ e~$m$ na fita, e pare com um único bloco composto
de $n+m$~$\TMstroke$. Os dois blocos de entrada de~$\TMstroke$ estão
separados por um~$\TMblank$, então um método seria escrever um traço no
quadrado que contém o $\TMblank$, e apagar o último~$\TMstroke$.
\begin{figure}\[
\begin{tikzpicture}[->,>=stealth',shorten >=1pt,auto,node distance=2.8cm,
                    semithick]
  \tikzstyle{every state}=[fill=none,draw=black,text=black]

  \node[initial,state]         (A)              {$q_0$};
  \node[state]         (B) [right of=A] {$q_1$};
  \node[state]         (C) [right of=B] {$q_2$};

  \path (A) edge node {\TMtrans{\TMblank}{\TMstroke}{\TMstay}} (B)
           edge [loop above] node {\TMtrans{\TMstroke}{\TMstroke}{\TMright}} (B)
        (B) edge [loop above] node {\TMtrans{\TMstroke}{\TMstroke}{\TMright}} (B)
            edge node {\TMtrans{\TMblank}{\TMblank}{\TMleft}} (C)
        (C) edge [loop above] node {\TMtrans{\TMstroke}{\TMblank}{\TMstay}} (C);
\end{tikzpicture}
\]\caption{Uma máquina que computa $f(x,y) = x+y$}
\ollabel{fig:adder}
\end{figure}
\end{ex}

\begin{prob}
Rastreie as configurações da máquina de
\olref[tur][mac][una]{ex:adder} para a entrada~$\tuple{3,2}$. O que acontece
se a máquina computa $0+0$?
\end{prob}

\begin{explain}
Em \olref[rep]{ex:doubler}, demos um exemplo de uma máquina de Turing
que recebe como entrada uma sequência de~$\TMstroke$ e para com uma sequência
de duas vezes mais~$\TMstroke$ na fita---a máquina duplicadora.
No entanto, como a saída contém $\TMblank$ à esquerda do
bloco duplicado de~$\TMstroke$, ela não computa de fato a
função $f(x) = 2x$, como você poderia ter suposto. Descreveremos duas
maneiras de corrigir isso.
\end{explain}

\begin{ex}
A máquina em \olref{fig:doubler-disc} computa a função $f(x) =
2x$. Em vez de apagar a entrada e escrever dois $\TMstroke$ na
extremidade direita para cada $\TMstroke$ da entrada, como faz a máquina de
\olref[rep]{ex:doubler}, esta máquina adiciona um único~$\TMstroke$
à direita para cada~$\TMstroke$ da entrada. Ela precisa rastrear
onde a entrada termina, então deixa um~$\TMblank$ entre a entrada
e os traços adicionados, que ela preenche com um~$\TMstroke$ bem ao
final. E temos que ``lembrar'' onde estamos na entrada, então
substituímos temporariamente um~$\TMstroke$ do bloco de entrada por um~$\TMblank$.
  \begin{figure}
\[
\begin{tikzpicture}[->,>=stealth',shorten >=1pt,auto,node distance=2.8cm,
                    semithick]
  \tikzstyle{every state}=[fill=none,draw=black,text=black]

  \node[initial,state] (0)              {$q_0$};
  \node[state]         (1) [above of=0] {$q_1$};
  \node[state]         (2) [above of=1] {$q_2$};
  \node[state]         (3) [right of=2] {$q_3$};
  \node[state]         (4) [below of=3] {$q_4$};
  \node[state]         (5) [below of=4] {$q_5$};
  \node[state]         (6) [above right of=4] {$q_6$};
  \node[state]         (7) [below of=6] {$q_7$};
  \node[state]         (8) [below of=7] {$q_8$};

  \path (0) edge node {\TMtrans{\TMstroke}{\TMblank}{\TMright}} (1)
%    (0) edge node {\TMtrans{\TMblank}{\TMblank}{\TMstay}} (h)
    (1) edge [loop left] node {\TMtrans{\TMstroke}{\TMstroke}{\TMright}} (1)
      edge node {\TMtrans{\TMblank}{\TMblank}{\TMright}} (2)
    (2) edge [loop above] node {\TMtrans{\TMstroke}{\TMstroke}{\TMright}} (2)
        edge node {\TMtrans{\TMblank}{\TMstroke}{\TMleft}} (3)
    (3) edge [loop above] node {\TMtrans{\TMstroke}{\TMstroke}{\TMleft}} (3)
        edge node[left] {\TMtrans{\TMblank}{\TMblank}{\TMleft}} (4)
    (4) edge        node[left] {\TMtrans{\TMstroke}{\TMstroke}{\TMleft}} (5)
    (5) edge [loop below] node {\TMtrans{\TMstroke}{\TMstroke}{\TMleft}} (5)
        edge              node {\TMtrans{\TMblank}{\TMstroke}{\TMright}} (0)
    (4) edge      node[sloped] {\TMtrans{\TMblank}{\TMstroke}{\TMright}} (6)
    (6) edge              node {\TMtrans{\TMblank}{\TMstroke}{\TMright}} (7)
    (7) edge [loop right] node {\TMtrans{\TMstroke}{\TMstroke}{\TMright}} (7)
        edge              node {\TMtrans{\TMblank}{\TMblank}{\TMleft}} (8)
    (8) edge [loop right] node {\TMtrans{\TMstroke}{\TMblank}{\TMstay}} (8);
    \end{tikzpicture}
\]
\caption{Uma máquina que computa $f(x) = 2x$}
\ollabel{fig:doubler-disc}
\end{figure}
\end{ex}

\begin{ex}\ollabel{ex:mover}
Uma segunda possibilidade para computar $f(x) = 2x$ é manter a máquina
duplicadora original, mas adicionar estados e instruções ao final que movam
o bloco duplicado de traços para a extremidade esquerda da fita.  A máquina
em \olref{fig:mover} faz exatamente esta última parte: iniciada em uma fita
composta de um bloco de~$\TMblank$ seguido de um bloco
de~$\TMstroke$ (e com a cabeça posicionada em qualquer lugar do bloco
de~$\TMblank$), ela apaga os $\TMstroke$ um a um e os escreve
no início da fita. Para conseguir saber quando ela
termina, primeiro marca o fim do bloco de $\TMstroke$ com
um símbolo~$\TMendtape$, que é apagado ao final. Começamos a
numerar os estados em~$q_6$, de modo que possam ser adicionados à máquina
duplicadora. Tudo de que você precisará é uma instrução adicional $\delta(q_0,
\TMblank) = \tuple{q_6,\TMblank,\TMstay}$, isto é, uma seta de~$q_0$
para~$q_6$ rotulada~$\TMtrans{\TMblank}{\TMblank}{\TMstay}$. (Há
um problema sutil: a máquina resultante não funciona para a
entrada~$x=0$. Deixaremos isso como exercício.)
\begin{figure}
  \[
  \begin{tikzpicture}[->,>=stealth',shorten >=1pt,auto,node distance=2.8cm,
                      semithick]
    \tikzstyle{every state}=[fill=none,draw=black,text=black]
    \node[initial,state] (6)              {$q_6$};
    \node[state]         (7) [right of=6] {$q_7$};
    \node[state]         (8) [right of=7] {$q_8$};
    \node[state]         (9) [below of=8] {$q_9$};
    \node[state]         (10) [left of=9] {$q_{10}$};
    \node[state]         (11) [left of=10]  {$q_{11}$};
    \node[state]         (12) [below of=11]  {$q_{12}$};
    \node[state]         (13) [right of=12] {$q_{13}$};
    \node[state]         (14) [right of=13] {$q_{14}$};
    \path
    (6)  edge node {\TMtrans{\TMblank}{\TMblank}{\TMright}} (7)
    (7)  edge [loop above] node {\TMtrans{\TMblank}{\TMblank}{\TMright}} (7)
         edge node {\TMtrans{\TMstroke}{\TMstroke}{\TMright}} (8)
    (8)  edge [loop above] node {\TMtrans{\TMstroke}{\TMstroke}{\TMright}} (8)
         edge node {\TMtrans{\TMblank}{\TMendtape}{\TMleft}} (9)
    (9)  edge [loop right] node {\TMtrans{\TMstroke}{\TMstroke}{\TMleft}} (9)
         edge node {\TMtrans{\TMblank}{\TMblank}{\TMright}} (10)
    (10) edge node[above] {\TMtrans{\TMstroke}{\TMblank}{\TMleft}} (11)
    (11) edge [loop above] node {\TMtrans{\TMblank}{\TMblank}{\TMleft}} (11)
         edge node[left] {\begin{tabular}{@{}l@{}}
          \TMtrans{\TMendtape}{\TMendtape}{\TMright}\\
          \TMtrans{\TMstroke}{\TMstroke}{\TMright}
         \end{tabular}} (12)
    (12) edge node[] {\TMtrans{\TMblank}{\TMstroke}{\TMright}} (13)
    (13) edge [loop below] node {\TMtrans{\TMblank}{\TMblank}{\TMright}} (13)
         edge node[sloped] {\TMtrans{\TMstroke}{\TMblank}{\TMleft}} (11)
    (13) edge node {\TMtrans{\TMendtape}{\TMblank}{\TMstay}} (14);
  \end{tikzpicture}
  \]
  \caption{Movendo um bloco de $\TMstroke$ para a esquerda}
  \ollabel{fig:mover}
\end{figure}
\end{ex}

\begin{prob}
Em \olref[tur][mac][una]{ex:mover} descrevemos uma máquina composta
de uma combinação da máquina duplicadora de
\olref[tur][mac][una]{fig:doubler-disc} e da máquina movedora de
\olref[tur][mac][una]{fig:mover}. O que acontece se você inicia essa
máquina combinada com a entrada~$x=0$, isto é, em uma fita vazia?  Como você
corrigiria a máquina de modo que, nesse caso, ela parasse com
saída~$2x=0$? (Você deve conseguir fazer isso adicionando um estado e
uma transição.)
\end{prob}

\begin{prob}
\emph{Subtração:} Projete uma máquina de Turing que, ao receber uma entrada
de duas cadeias não vazias de traços de comprimento $n$ e~$m$, em que $n >
m$, compute a função $f(n,m) = n - m$.
\end{prob}

\begin{prob}
\emph{Igualdade:} Projete uma máquina de Turing para computar a seguinte função:
\[
\fn{igualdade}(n,m) = 
\begin{cases}
  \text{1} & \text{se~$n = m$} \\
  \text{0} & \text{se~$n \neq m$}
\end{cases}
\]
em que~$n$ e~$m \in \PosInt$.
\end{prob}

\begin{prob}
Projete uma máquina de Turing para computar a função $\min(x,y)$, em que $x$
e $y$ são inteiros positivos representados na fita por cadeias de
$\TMstroke$ separadas por um $\TMblank$. Você pode usar símbolos
adicionais no alfabeto da máquina.

A função $\min$ seleciona o menor valor entre seus argumentos, de modo que
$\min(3,5)=3$, $\min(20,16)=16$, e $\min(4,4)=4$, e assim por diante.
\end{prob}

\begin{defn}
  Uma máquina de Turing~$M$ computa a função parcial $f\colon \Nat^k
  \pto \Nat$ sse, 
  \begin{enumerate}
    \item $M$ para com a entrada $\TMstroke^{n_1}\concat\TMblank\concat
    \dots \concat\TMblank\concat\TMstroke^{n_k}$ produzindo a saída $\TMstroke^{m}$ se $f(n_1, \dots, n_k) = m$.
    \item $M$ não para de modo algum, ou para com uma saída que não é um
    único bloco de~$\TMstroke$ se $f(n_1, \dots, n_k)$ é indefinida.
  \end{enumerate}
\end{defn}

\end{document}
